\scanpage{150}
Ceci montre que l'on est réduit à étudier \inserted{$\psi_\Sigma$ dans} le cas où
$K=K_0$, $E=E_0$ (donc $E,K$ weiliens et finis).

Identifiant les \unsure{car.} de $T_E,T_K$ à valeurs dans $\overline{\mathbb Q}$ à
$\mathbb Z^{I_E}$, $\mathbb Z^{I_K}$, on constate que l'hom
\inserted{(\unsure{compatible à $G$})} sur les \deleted{\unsure{car.}} groupes de caractères
défini par $\unsure{T_E|\psi_\Sigma}$:
\diagram{\begin{tikzcd}\mathbb Z^{I_E}\arrow[r,no head]&\mathbb Z^{I_K}\end{tikzcd}}
est donné par la matrice $1_\Sigma$ (considéré comme une fonction
\deleted{caractéristique} sur $I_E\times I_K$).
\uline{Donc la connaissance de} $\Sigma$ \uline{équivaut à celle de} $\psi_\Sigma$
(cette dernière n'étant d'ailleurs pas quelconque\ldots).

Choisissons une clôture algébrique $\overline E$ de $E$,
\deleted{\unsure{supposons $K$}} \inserted{supposons $K$ fini (degré $2d'$) sur $\mathbb Q$,}
alors la structure $\Sigma$ équivaut à la donnée d'un ensemble \deleted{de} d'hom
\[
\sigma_1,\ldots\sigma_{d'}:K\longrightarrow\overline E
\]
\deleted{\unsure{conjugués}} tels que pour tte conj. $\tau$ de $G=\operatorname{Gal}(\overline E/\mathbb Q)$,
les $\{\tau\sigma_i,\sigma_i\}$ font l'ens de tous les hom $K\longrightarrow\overline E$.
On aura alors
\[
\psi_\Sigma(x)=\sigma_1(x)\sigma_2(x)\cdots\sigma_{d'}(x)
\]
Si $K$ est weilien,
\deleted{\unsure{donc} si $x\in K_r$ (ss-corps tot. réel),}
\inserted{les $\sigma_i|K_r$ font \deleted{\unsure{un ens}} l'ens des hom
$K_r-\overline E$, et on trouve}
\deleted{on aura $\tau\sigma_i(x)=\sigma_i(x)$, donc}
\[
\deleted{\psi_\Sigma(x)^2=\unsure{N_K}}
\]
\[
\psi_\Sigma|T_{K_r}=N_{K_r/\mathbb Q}
\]
Par suite, $\psi_\Sigma$ s'annule sur $V_{K_r}\subset T_{K_r}\subset T_K$.
De ceci, et du §1, on conclut que \uline{si} $K$ \uline{fini sur} $\mathbb Q$,
alors $\psi_\Sigma:T_K\longrightarrow T_E$ \uline{se factorise par} $S_K^0\longrightarrow T_E$:
\margin{149}
